% Complete local avoidance constructions. All mathematical statements
% are independent of the source-note numbering.
\section{Local avoidance constructions}\label{app:local}
Throughout this appendix, the family is $k$-uniform and has transversal
number greater than the integer request budget $T\le k$. Thus every set
of at most $T$ points is avoided by a family edge. All edges obtained by
such requests belong to the original family. Repetitions cause no problem:
the contradiction is a subfamily of at most seven edges.

The notation of a good triple is as in Section~\ref{sec:prelim}.
A construction \emph{closes} a triple if four further requests force a
subfamily with no transversal of size at most two. Pair-cell sizes in
this appendix are unnormalised integers. Whenever gap parameters and a
real budget coefficient occur, assume $0\le\ell\le h\le1$ and
$0\le\beta\le1$.

\begin{lemma}[Three small intersections; cf.~\cite{FKW}]\label{lem:smalltriple}
Suppose all three pairwise intersections of a good triple have size at most $m\le k/2$. It extends to a bad subfamily of at most seven edges whenever
\[
\tau>B:=\left\lceil\max\{(3k+m)/4,(2k+2m)/3\}\right\rceil.
\]
\end{lemma}
\begin{proof}
Write the triple as $E,F,G$, with $X=E\cap G$, $Y=E\cap F$, $Z=F\cap G$, and relabel so $x=|X|\ge y=|Y|\ge z=|Z|$. Choose $P_G\subseteq G\setminus(E\cup F)$ of size $B-x-y$, and $B_0\subseteq F\setminus(E\cup G)$ of size $\min(B-x,k-y-z)$. Put $C=F\setminus(Y\cup B_0)$; then
\[
c=|C|=\max(k-B+x-y,z).
\]
Choose $P_E\subseteq E\setminus(F\cup G)$ of size $\min(B-x-c,k-x-y)$. All requested sizes are nonnegative and fit their hosts: $B-x-y\ge B-2m\ge0$, $c\le k-B+m$, $B-x-c\ge2B-k-2m\ge0$, and $B-x-y\le k-x-z$ because $y\ge z$ and $B\le k$.

Request edges avoiding
\[
X\cup Y\cup P_G,\quad X\cup B_0,\quad X\cup C\cup P_E,\quad (E\cup G)\setminus(X\cup P_E\cup P_G).
\]
The first three sizes are at most $B$. The fourth size is
\[
\max\{3(k-B)+x,\ 2(k-B)+y+z,\ (k-B)+2y\}\le B,
\]
by the definition of $B$.
A candidate pair with one point in $X$ has its other point in $Y,B_0$, or $C$ and is missed by one of the first three responses. Otherwise a point in $Y$ must pair with a point in $G\setminus X$, and the first or fourth response misses the pair according as the latter point is in $P_G$ or its complement. If the first point is in $Z$, use the third or fourth response and $P_E$. This exhausts the candidate pairs.
\end{proof}

\begin{lemma}[Splitting one pair cell]\label{lem:split}
For a good triple with pair sizes $x=|E\cap F|$, $y=|E\cap G|$, $z=|F\cap G|$, four legal responses yield a bad seven-tuple at any budget $T\le k$ satisfying
\[
T\ge\max\{(k+x+y+z)/2,\ k-x+|y-z|,\ k/3+x\}.
\]
The lemma is exactly integral.
\end{lemma}
\begin{proof}
Put $S=x+y+z$; the first hypothesis and $T\le k$ give $S\le T$. Request $H$ avoiding $X=E\cap F$, $Y=E\cap G$, $Z=F\cap G$, and
\[
p=\max(0,k-2(T-x)-y-z)
\]
points of the private part of $G$. This selection fits because $T\ge x$. Its union size is
\[
S+p=x+\max(y+z,k-2T+2x)\le T.
\]
All four triple cells vanish. Put $A=F\cap H$, $B=G\cap H$, $C=E\cap H$ with sizes $a,b,c$. Then
\[
a\le k-x-z,\quad b\le2(T-x),\quad c\le k-x-y,\quad a+b+c\le k.
\]

If $b\le T-x$, the three requests $X\cup B$, $Y\cup A$, $Z\cup C$ have sizes at most $T$: for the latter two use $k-x+|y-z|\le T$. They cover the three complementary candidate pairs.

If $b>T-x$, divide $B$ into two parts of size at most $T-x$; this is possible because $b\le2(T-x)$, also integrally. Two requests contain $X$ and one part each. The third contains $Y\cup Z\cup A\cup C$, whose size satisfies
\[
y+z+a+c\le k-b+y+z<k-T+x+y+z\le T.
\]
These requests again cover all three candidate pairs.
\end{proof}

\begin{lemma}[One surviving triple cell]\label{lem:fourcase}
Let $E,F,G$ be a good triple of $k$-sets, and put $X=E\cap F$, $Y=E\cap G$, $Z=F\cap G$, with sizes $x,y,z$ and $S=x+y+z$. Four further avoidance responses force a bad seven-tuple at any integer budget $T\le k$ satisfying
\[
\begin{gathered}
T\ge\max\Bigl\{x+y,\ k/2+x,\ k/2+y,\ k+x-y-z,\ k-x+y-z,\\
 k-S/3,\ (3k+S)/5,\ (k+x+y+2z)/3,\ (2k+3z)/4\Bigr\}.
\end{gathered}
\]
The construction is exactly integral. Any of the three pair cells can be designated $Z$.
\end{lemma}
\begin{proof}
Request $H$ avoiding $X,Y$ and a subset $Z_0\subseteq Z$ of size $\min(z,T-x-y)$. This is legal. Write
\[
Q=Z\cap H,\quad A=(F\cap H)\setminus Q,\quad B=(G\cap H)\setminus Q,\quad C=E\cap H,
\]
with sizes $q,a,b,c$. Also put $V=Z\setminus Q$, $v=z-q$, and $D=E\setminus(X\cup Y\cup C)$, $d=k-x-y-c$. These eight cells are disjoint, and
\[
q\le(S-T)_+,\quad a\le k-x-z,\quad b\le k-y-z,\quad c\le k-x-y,\quad q+a+b+c\le k.
\]
The only triple cell of $E,F,G,H$ is $Q$. A pair piercing these four edges therefore either belongs to $Q\times E$, or to one of $X\times B$, $Y\times A$, $V\times C$. We will make all three remaining avoidance requests contain $Q$ and their union contain $E$, while covering these last three products.

Two uniform bounds used below are
\[
q+x+b\le\max(k+x-y-z,k+2x-T)\le T,
\]
\[
q+y+a\le\max(k-x+y-z,k+2y-T)\le T.
\]
Also $z\le T$, since $4T\ge2k+3z$ and $z\le k$.

\paragraph{Case 0: $c\le T-z$.} Start the three requests with the disjoint-cell unions
\[
Q\cup X\cup B,\qquad Q\cup Y\cup A,\qquad Z\cup C.
\]
Each base has size at most $T$. Their total load, plus $d$, is
\[
k+2q+a+b+z\le\max(3k-S,3k+S-2T)\le3T.
\]
The first inequality follows from $a+b\le2k-x-y-2z$ and $q\le(S-T)_+$. Distribute $D$ among the remaining capacities. These capacities and $d$ are integers; hence the distribution can be integral. All candidate pairs are now covered by the bases or by $Q$ together with this distribution.

In the other cases $c>T-z$. Set
\[
a_0=k+x+z-2T,\qquad b_0=k+y+z-2T.
\]

\textbf{Case 1: $a<a_0$.} Use bases
\[
Q\cup X\cup B,\qquad Y\cup A\cup Z\cup C_2,\qquad Z\cup C_3,
\]
where $C=C_2\sqcup C_3$. This split fits because
\[
y+a+z\le k+x+y+2z-2T\le T,
\]
\[
a+c\le2k+z-y-2T\le2T-2z-y.
\]
Thus the two capacities $T-y-a-z$ and $T-z$ are nonnegative and sum to at least $c$. The total base load plus $d$ is
\[
k+q+a+b+2z\le2k+2z-c<2k+3z-T\le3T.
\]
Distribute $D$ among the remaining capacities. The three products are covered, and all bases contain $Q$.

\textbf{Case 2: $b<b_0$.} Interchange $X,B,x,b$ with $Y,A,y,a$ in Case 1. This covers the case even if both inequalities hold.

\textbf{Case 3: $a\ge a_0$ and $b\ge b_0$.} Put $u=c+z-T>0$. Since $z\le T$, we have $u\le c$, so choose $C_{12}\subseteq C$ of size $u$ and put $C_3=C\setminus C_{12}$. Choose an integer $t$ in
\[
\max(0,y+a+z+u-T)\ \le t\le\ \min(v,T-q-x-b-u).
\]
This interval is nonempty. The second upper entry is nonnegative because
\[
q+b+c\le k-a\le2T-x-z.
\]
The second lower entry is at most $v$ because
\[
q+a+c\le k-b\le2T-y-z.
\]
Finally the untruncated lower entry does not exceed the untruncated upper entry: this is equivalent to
\[
q+a+b+2c+x+y+3z\le4T,
\]
which follows from $q+a+b+c\le k$, $c\le k-x-y$ and $2k+3z\le4T$. All endpoints are integers.

Split $V=V_{13}\sqcup V_{23}$ with $|V_{13}|=t$, and use bases
\[
Q\cup X\cup B\cup C_{12}\cup V_{13},
\]
\[
Q\cup Y\cup A\cup C_{12}\cup V_{23},
\]
\[
Z\cup C_3.
\]
The interval inequalities make the first two sizes at most $T$; the third has size exactly $T$. Their total load plus $d$ is
\[
k+q+a+b+2z+u\le2k+3z-T\le3T.
\]
Again distribute $D$ into the residual capacities. The first two requests cover $X\times B$, $Y\times A$ and $V\times C_{12}$; the third covers $V\times C_3$. All contain $Q$ and their union contains $E$. Thus each candidate piercing pair of the first four edges misses at least one of the three new responses. This is a bad seven-tuple.
\end{proof}

\begin{lemma}[An asymmetric split]\label{lem:asymone}
With the good-triple notation of Lemma~\ref{lem:fourcase}, four further avoidance responses force a bad seven-tuple whenever $T\le k$ and
\[
T\ge\max\left\{S,\ \frac k2+y,\ \frac{k+2x-y+z}{2},\ \frac{k+2x+y+3z}{3}\right\}.
\]
The construction is exactly integral. All six permutations of $x,y,z$ are available.
\end{lemma}
\begin{proof}
Request $H$ avoiding $X,Y,Z$ and $T-S$ points in the private part of $F$. The private selection fits because $T-S\ge0$ and $T-S\le k-x-z$, the latter following from $T\le k+y$. Put $A=F\cap H$, $B=G\cap H$, $C=E\cap H$, with sizes $a,b,c$. All four triple cells vanish, so the candidate piercing pairs are $X\times B$, $Y\times A$, and $Z\times C$. Also
\[
a\le k-T+y,\qquad b\le k-y-z,\qquad a+b+c\le k.
\]
If $a\le T-y-z$, split $B=B_1\sqcup B_2$ into two parts of size at most $T-x-z$. This is possible, integrally, since
\[
2(T-x-z)\ge k-y-z\ge b.
\]
Start the final three requests with
\[
X\cup Z\cup B_1,\qquad X\cup Z\cup B_2,\qquad Y\cup Z\cup A.
\]
The first two fit by construction; the third fits by the case assumption. The total base load plus $c$ is
\[
2x+y+3z+a+b+c\le k+2x+y+3z\le3T.
\]
Distribute $C$ integrally among the residual capacities. Every final request contains $Z$, so this distribution covers $Z\times C$; the first two requests cover $X\times B$ and the third covers $Y\times A$.

If instead $a>T-y-z$, then
\[
b+c\le k-a<k-T+y+z\le2(T-x-z).
\]
Split $B\cup C$ into two parts of size at most $T-x-z$, and put $X\cup Z$ together with one part in each of the first two requests. The third request is $Y\cup A$, of size at most $k-T+2y\le T$. These three requests cover all three candidate-pair products again. Thus no pair pierces the seven edges in either case.
\end{proof}

\begin{lemma}[A second asymmetric split]\label{lem:asymtwo}
With the same good-triple notation, four further avoidance responses force a bad seven-tuple if $T\le k$ and
\[
T\ge\max\left\{S,\ \frac k3+x,\ k-x+z,\ \frac{k+2x+3y+z}{3}\right\}.
\]
The construction is exactly integral, with all six permutations available.
\end{lemma}
\begin{proof}
Request $H$ avoiding $X,Y,Z$ and
\[
p=\max(0,k-2(T-x)-y-z)
\]
points of the private part of $G$. This fits its host since $T\ge x$, and its total cost is $\max(S,k+3x-2T)\le T$. As in Lemma~\ref{lem:split}, write $A=F\cap H$, $B=G\cap H$, $C=E\cap H$, with sizes $a,b,c$. Then
\[
b\le2(T-x),\qquad c\le k-x-y,\qquad a+b+c\le k.
\]
All four triple cells vanish and the candidate pairs are $X\times B$, $Y\times A$, $Z\times C$.

If $b>k+y+z-T$, split $B$ into two parts of size at most $T-x$, and use $X$ with each part as two requests. The third is $Y\cup Z\cup A\cup C$, of size at most $k-b+y+z<T$.

Otherwise $b\le k+y+z-T\le2(T-x-y)$, where the last inequality is exactly $3T\ge k+2x+3y+z$. Split $B=B_1\sqcup B_2$ into two parts of size at most $T-x-y$. Start the three requests with
\[
X\cup Y\cup B_1,\qquad X\cup Y\cup B_2,\qquad Y\cup Z\cup C.
\]
The first two fit by construction; the third has size at most $k-x+z\le T$. The total base load plus $a$ is at most
\[
2x+3y+z+a+b+c\le k+2x+3y+z\le3T.
\]
Distribute $A$ into the residual integer capacities. Since every request contains $Y$, this covers $Y\times A$. The first two cover $X\times B$ and the third covers $Z\times C$. Both cases therefore yield a bad seven-tuple.
\end{proof}

\begin{lemma}[Two traces forced through a gap]\label{lem:gapempty}
Suppose no pair of family edges has intersection in $[\ell k,hk]$. For a good triple with pair sizes $x,y,z$, define $S=x+y+z$. If each of
\[
S,\quad (1-h)k+y,\quad (1-h)k+z,\quad (2-2h)k-x,
\]
\[
k/3+x,\quad ((2-h)k+2x-y)/3,\quad ((2-h)k+2x-z)/3,
\]
\[
((3-2h)k+x-y-z)/3,\quad 2\ell k+y+z
\]
is at most $\beta k$, and $T=\lceil\beta k\rceil+1\le k$, four further requests of size at most $T$ produce a bad seven-tuple. Here $x$ is the pair cell shared by the two edges whose new traces will be forced below $\ell k$.
\end{lemma}
\begin{proof}
Write the edges as $E,F,G$, with $X=E\cap F$, $Y=E\cap G$, $Z=F\cap G$. Let $h_0=\lfloor hk\rfloor$ and choose private subsets of $E,F$ of sizes
\[
p=(k-x-y-h_0)_+,\qquad t=(k-x-z-h_0)_+.
\]
Avoid them together with $X,Y,Z$. The cost before padding is
\[
C_0=x+\max(y,k-x-h_0)+\max(z,k-x-h_0).
\]
Expanding these two maxima gives the first four displayed forms, with rounding error strictly less than two. Thus $C_0<\beta k+2$, and its integrality gives $C_0\le\lceil\beta k\rceil+1=T$. Pad the request with $T-C_0$ private points of $G$. This fits: $C_0\ge x+y+z$ and $T\le k$ imply $T-C_0\le k-y-z$.

The response $H$ has its $E$- and $F$-traces at most $h_0$, so the gap makes both strictly smaller than $\ell k$. Its $G$-trace has size $b\le k-T+x+p+t$. All triple intersections among $E,F,G,H$ are empty. Split the $G$-trace into two nearly equal integer parts and use $X$ together with each part in two requests. Their sizes are at most
\[
\frac{k+3x+p+t-T}{2}+\frac12.
\]
Expanding the two positive parts in $p,t$, the numerator $k+3x+p+t$ is the maximum of
\[
k+3x,\quad (2-h)k+2x-y,\quad (2-h)k+2x-z,\quad (3-2h)k+x-y-z,
\]
with error strictly less than two. The next four hypotheses therefore make each split request strictly less than $(3\beta k+2-T)/2+1/2\le T$, since $T\ge\beta k+1$.

Use $Y,Z$ and the two small traces of $H$ in the third request. Its size is strictly smaller than $y+z+2\ell k\le\beta k$. These three requests eliminate the three complementary candidate-pair products.
\end{proof}

\begin{lemma}[A surviving triple cell and a gap]\label{lem:gapsurvive}
Suppose no pair of family edges has intersection in $[\ell k,hk]$. For a good triple with the usual pair sizes $x,y,z$, put $S=x+y+z$. Suppose $S\ge\beta k$, and each of
\[
x+y,\quad (1-h)k+y,\quad (1-h)k+z,
\]
\[
k/4+x+(y+z)/4,\quad 2\ell k+y+z,\quad k/2+x/2+(z+\ell k)/4
\]
is at most $\beta k$. If $T=\lceil\beta k\rceil+1\le k$, four further requests of size at most $T$ force a bad seven-tuple. The assertion includes the rounding strip $\beta k\le S<T$.
\end{lemma}
\begin{proof}
Request $H$ avoiding $X,Y$ and a subset $Z_0\subseteq Z$ of size $\min(z,T-x-y)$. This fits its budget since $x+y\le\beta k\le T$. Put $Q=Z\cap H$, $A=(F\cap H)\setminus Q$, $B=(G\cap H)\setminus Q$, $C=E\cap H$, with sizes $q,a,b,c$. Then
\[
q\le(S-T)_+\le S-\beta k,\qquad b\le k-y-z.
\]
Moreover
\[
c\le k-x-y=k-S+z\le(1-\beta)k+z\le hk,
\]
\[
a+q\le k-x-z+(S-T)_+\le(1-\beta)k+y\le hk.
\]
For the second inequality, use $S\ge\beta k$ if $S\le T$, and $T\ge\beta k$ otherwise. These are actual pair-intersection sizes, so the gap forces $c<\ell k$ and $a+q<\ell k$.

Split $B=B_1\sqcup B_2$ into nearly equal integer parts. Start two requests with $X\cup Q\cup B_1$ and $X\cup Q\cup B_2$. Their sizes are at most
\[
x+q+\lceil b/2\rceil\le2x+(y+z)/2+k/2-\beta k+1/2\le\beta k+1/2\le T.
\]
Start the third request with $Y\cup Z\cup A\cup C$, of size strictly below $y+z+2\ell k-q\le\beta k$.

All three bases contain $Q$. To make their union contain all of $E$, distribute $D=E\setminus(X\cup Y\cup C)$ among the residual capacities. The total base load plus $|D|$ is
\[
k+x+z+2q+a+b<k+x+z+q+\ell k+b
\]
\[
\le2k+x-y+\ell k+q\le2k+2x+z+\ell k-\beta k\le3\beta k\le3T.
\]
The distribution therefore exists with integer sizes. Candidate piercing pairs of $E,F,G,H$ are in $Q\times E$, $X\times B$, $Y\times A$, or $(Z\setminus Q)\times C$. The first kind is eliminated because every final request contains $Q$ and their union contains $E$; the other three kinds are eliminated by the indicated bases.
\end{proof}

\begin{lemma}[Two large opposite pair cells]\label{lem:twolarge}
Suppose all pair intersections in the family are at most $m$ or greater than $k/2$, where $m\le k/4$. Let $E,F,G,H$ be four $k$-edges with all triple intersections empty. Put $X=E\cap F$, $B=G\cap H$, of sizes $x,b>k/2$, and suppose the other four pair cells each have size at most $m$. Three further legal responses give a bad seven-tuple if $T\le k$ and
\[
T\ge\lceil k/2\rceil+2m,\qquad T\ge x+m,\qquad 3T\ge2x+b+\lceil k/2\rceil+2m.
\]
\end{lemma}
\begin{proof}
Put $Y=E\cap G$, $Z=F\cap G$, $A=F\cap H$, $C=E\cap H$, and let $U=Y\cup Z\cup A\cup C$. These six pair cells are disjoint. Write $s=|Y|+|Z|$, $t=|A|+|C|$, so $s,t\le2m$, and define
\[
p=\lceil k/2\rceil-\min(s,t),\qquad q=T-\lceil k/2\rceil-\max(s,t).
\]
Both are nonnegative. Moreover $p\le\lceil k/2\rceil\le b$ and $q\le T-\lceil k/2\rceil\le\lfloor k/2\rfloor<x$. Choose $B_0\subseteq B$, $X_0\subseteq X$ with sizes $p,q$, and request $I$ avoiding $U\cup B_0\cup X_0$, whose size is exactly $T$.

The $G$-trace of $I$ has size at most
\[
k-s-p=\lfloor k/2\rfloor+\min(s,t)-s\le\lfloor k/2\rfloor,
\]
and the analogous statement holds for $H$. The global gap therefore forces both traces to have size at most $m$. In particular $d=|B\cap I|\le m$. Also $x_I=|X\cap I|\le x-q$.

Choose $B_1\subseteq B$ containing $B\cap I$ and having size $\min(b,T-x)$. This is possible since $T-x\ge m\ge d$. Use the last two requests
\[
X\cup B_1,\qquad (X\cap I)\cup(B\setminus B_1).
\]
The first has size at most $T$. For the second, note that
\[
x+x_I+b\le2x+b-q\le2x+b-T+\lceil k/2\rceil+2m\le2T.
\]
It follows that its size $x_I+\max(0,b-T+x)$ is at most $T$ as well.

A pair piercing $E,F,G,H$ belongs to $X\times B$, $Y\times A$ or $Z\times C$. The latter two products lie entirely in $U$ and hence miss $I$. For a remaining pair in $X\times B$, if its $B$-point belongs to $B_1$, the penultimate response misses both points. Otherwise its $B$-point lies outside $I$, since $B\cap I\subseteq B_1$; its $X$-point must then belong to $I$, and the last response misses both. This proves the lemma.
\end{proof}

\begin{lemma}[Two surviving triple cells]\label{lem:gaptwo}
Suppose no pair intersection lies in $[\ell k,hk]$. For a good triple $E,F,G$ let $X=E\cap F$, $Y=E\cap G$, $Z=F\cap G$, with sizes $x,y,z$, and put $S=x+y+z$. Suppose
\[
x+y\ge(1-h)k,\qquad y+z\ge(1-h)k,
\]
and each of
\[
(2-2h)k-y,\quad x+\ell k,\quad z+\ell k,\quad k/2+y,\quad 3k/4,\quad (3k+S)/5
\]
is at most $\beta k$. Then four further responses of avoidance budget $T=\lceil\beta k\rceil+1\le k$ give a bad seven-tuple. All six orientations are allowed.
\end{lemma}
\begin{proof}
Let $h_0=\lfloor hk\rfloor$ and $g=(k-y-h_0)_+$. The domain inequalities imply $g\le x,z$. Choose $g$ points from each of $X,Z$, and avoid them together with all of $Y$. The initial cost is
\[
y+2g=\max(y,2k-y-2h_0)<\beta k+2.
\]
Since the cost is an integer, it is at most $\lceil\beta k\rceil+1=T$. Complete the portion of the request inside $F$ to exactly $T-y$ points, first using any remaining points of $X\cup Z$ and then private points of $F$. This is possible because $X,Z$ are disjoint subsets of $F$, $Y$ is disjoint from $F$, and $T-y\le k$. Let $H$ be an avoiding response.

Write $P=E\cap F\cap H$ and $Q=F\cap G\cap H$, of sizes $p,q$. No point of $Y$ is in $H$, so these are the only possible triple cells among the four edges. Set
\[
A=(F\cap H)\setminus(P\cup Q),\quad B=(G\cap H)\setminus Q,\quad C=(E\cap H)\setminus P,
\]
with sizes $a,b,c$. The priority given to $X\cup Z$ in the first request yields
\[
p+q\le(S-T)_+,
\]
and its total size inside $F$ gives
\[
d:=|F\cap H|=p+q+a\le k-T+y.
\]
The two $g$-point cuts also give $|E\cap H|,|G\cap H|\le h_0$. The gap therefore makes both intersections smaller than $\ell k$.

Start the final three requests with
\[
X\cup(G\cap H),\qquad Z\cup(E\cap H),\qquad Y\cup(F\cap H).
\]
Their sizes are at most $x+\ell k$, $z+\ell k$, and $k-T+2y$, respectively, so all fit in $T$. Each contains both $P$ and $Q$.

Let $D=E\setminus(F\cup G\cup H)$ and $W=G\setminus(E\cup F\cup H)$. These are disjoint, with sizes $k-x-y-c$ and $k-y-z-b$. The sum of the three base sizes plus $|D|+|W|$ is
\[
2k-y+2(p+q)+a=2k-y+(p+q)+d\le3k-T+(S-T)_+.
\]
This is at most $3T$: if $S\le T$, use $T\ge3k/4$; otherwise use $5T\ge3k+S$. Consequently $D\cup W$ can be partitioned among the residual integer capacities of the three requests.

Their union now contains $E\cup G$. A candidate piercing pair involving $P$ has its other point in $G$, while a pair involving $Q$ has its other point in $E$; all such pairs are eliminated since every request contains $P,Q$ and their union contains $E,G$. Every other candidate pair lies in $(X\setminus P)\times B$, $(Z\setminus Q)\times C$, or $Y\times A$, contained in the first, second, or third request respectively. Their three avoiding responses give the contradiction.
\end{proof}
